\section{Classical formulas used in the integral representation}\label{app:classical}
\source{Euler's formula, Hermite's formula, and the partial fractions of
\(\coth\) are classical; exact references are \cite{dlmf}.
This appendix expands the derivation in the earlier draft of this chat.}

The moment and pole calculations use two familiar facts: Euler's formula
for even zeta values and a special Hurwitz zeta integral. Both follow from
one elementary identity. Put $f(y)=(e^{2\pi y}-1)^{-1}$. For $u>0$,
\begin{equation}\label{eq:kernel}
 \frac1{1-e^{-u}}=\frac12+\frac1u
       +2\int_0^\infty f(y)\sin(uy)\dd y.
\end{equation}

To prove it, periodically extend
$F_u(x)=\cosh(u(x-\tfrac12))$ from $[0,1]$. Its endpoints agree,
and its Fourier coefficients satisfy
\[
 c_n=\int_0^1 F_u(x)e^{-2\pi\mathrm i nx}\dd x
     =\frac{2u\sinh(u/2)}{u^2+4\pi^2n^2},\qquad n\in\Z.
\]
Indeed, integrate $F_u''=u^2F_u$ twice by parts. The equal endpoint
values cancel, and the endpoint derivative difference is
$2u\sinh(u/2)$, giving
$(u^2+4\pi^2n^2)c_n=2u\sinh(u/2)$.
The periodic function is continuous and piecewise smooth. The Fourier
convergence theorem at $x=0$ therefore gives
$\cosh(u/2)=c_0+2\sum_{n\geq1}c_n$. Dividing by $2\sinh(u/2)$ yields
\[
 \frac12\coth(u/2)=\frac1u
      +2u\sum_{n\geq1}\frac1{u^2+4\pi^2n^2}.
\]
Expanding $f$ as a geometric series and integrating each exponential gives
\[
 \int_0^\infty f(y)\sin(uy)\dd y
 =\sum_{n\geq1}\int_0^\infty e^{-2\pi ny}\sin(uy)\dd y
 =\sum_{n\geq1}\frac{u}{u^2+4\pi^2n^2}.
\]
This interchange is justified by $|\sin(uy)|\leq uy$:
the sum of the absolute integrals is at most
$u\sum_{n\geq1}(2\pi n)^{-2}<\infty$.
Together with $1/(1-e^{-u})=\tfrac12+\tfrac12\coth(u/2)$,
this proves \eqref{eq:kernel}.

\textbf{Euler's formula.}
The same partial fractions give, for $0<u<2\pi$,
\begin{align*}
 \frac{u}{e^u-1}
 &=1-\frac u2+2u^2\sum_{n\geq1}\frac1{u^2+4\pi^2n^2}\\
 &=1-\frac u2+
   2\sum_{r\geq0}\frac{(-1)^r\zeta(2r+2)}{(2\pi)^{2r+2}}u^{2r+2}.
\end{align*}
For the second equality, expand each denominator geometrically. The
resulting double series converges absolutely whenever $|u|<2\pi$, so
the order of summation may be exchanged. Compare with the defining
series $u/(e^u-1)=\sum_{k\geq0}B_ku^k/k!$. Equality of coefficients gives
\[
 \boxed{\zeta(2e+2)=
 \frac{(-1)^eB_{2e+2}(2\pi)^{2e+2}}{2(2e+2)!}},
 \qquad e\in\Z_{\geq0}.
\]
This is the identity used to turn the polynomial integrals into the
rational Bernoulli values defining $\mu$.

\subsection{The Hurwitz integral}

Fix $a>0$. The definition
$\zeta(5,a)=\sum_{n\geq0}(a+n)^{-5}$ and the elementary exponential
integral give
\[
 \frac1{(a+n)^5}=\frac1{24}
       \int_0^\infty u^4e^{-(a+n)u}\dd u.
\]
The integrands are nonnegative, so Tonelli's theorem allows the sum and
integral to be exchanged. Summing the geometric series gives
\[
 \zeta(5,a)=\frac1{24}
       \int_0^\infty\frac{u^4e^{-au}}{1-e^{-u}}\dd u.
\]
Insert \eqref{eq:kernel}. The first two terms can be evaluated immediately:
\[
 \frac1{24}\int_0^\infty u^4e^{-au}
       \left(\frac12+\frac1u\right)\dd u
       =\frac1{2a^5}+\frac1{4a^4}.
\]
The remaining term is a double integral. Although the sine can change
sign, Fubini's theorem applies because
\begin{align*}
 &\int_0^\infty\!\int_0^\infty
       u^4e^{-au}f(y)|\sin(uy)|\dd y\dd u\\
 &\hspace{12mm}\leq
       \left(\int_0^\infty u^5e^{-au}\dd u\right)
       \left(\int_0^\infty yf(y)\dd y\right)<\infty.
\end{align*}
The first factor is $5!/a^6$. For the second, $yf(y)$ tends to
$1/(2\pi)$ at zero and decays exponentially at infinity. Thus the bound
also justifies subtracting and rearranging the integrals in the preceding
calculation.

After exchanging the integrations, evaluate the inner integral by taking
the imaginary part of an exponential integral:
\[
 \int_0^\infty u^4e^{-au}\sin(uy)\dd u
 =\operatorname{Im}\int_0^\infty
       u^4e^{-(a-\mathrm i y)u}\dd u
 =24\operatorname{Im}(a-\mathrm i y)^{-5}.
\]
The last equality follows by four integrations by parts; every boundary
term vanishes because the real part of $a-\mathrm i y$ is positive.
Keeping the factor $2/24$ from \eqref{eq:kernel}, we obtain
\[
 \boxed{\zeta(5,a)=\frac1{2a^5}+\frac1{4a^4}
 +2\int_0^\infty
   \frac{\operatorname{Im}(a-\mathrm i y)^{-5}}{e^{2\pi y}-1}\dd y.}
\]
This is precisely the special Hermite formula \eqref{eq:hermite} used in
the pole calculation. In particular, it holds for every positive integer
$a=j$, where the Hurwitz zeta value is the tail of the series for
$\zeta(5)$.

The two classical identities have now been derived from the same kernel.
The remaining passage to the positive weight $w$ is the integration by
parts carried out in the main text. Together, these calculations justify
the integral representation on both parts of the partial-fraction basis.

\bigskip
{\small\textbf{Classical references.}
The identities derived here are also recorded in the NIST Digital Library
of Mathematical Functions:
\href{https://dlmf.nist.gov/4.36.E3}{the partial fractions for $\coth$ (4.36.3)},
\href{https://dlmf.nist.gov/25.6.E2}{Euler's even-zeta formula (25.6.2)}, and
\href{https://dlmf.nist.gov/25.11.E29}{Hermite's Hurwitz integral (25.11.29)}.}
